% 第九卷 12.7.1.1：教学模型上的网格束搜索（每格宽度k=1，两条双词元约束ab、cd，C=4）。
% 格中数值为累计概率，由正文给出的教学分布逐步相乘得到；第5步只画出交付的假设。
\begin{tikzpicture}[
  x=0.9mm,y=1mm,figure picture,
  cell/.style={rectangle,rounded corners=1mm,draw=FigureStroke,fill=white,
    line width=.6pt,minimum width=19mm,minimum height=7mm,inner sep=1pt,
    font=\figurefont\scriptsize,align=center},
  best cell/.style={cell,draw=BookTeal,fill=FigureModelFill,line width=1pt},
  done cell/.style={cell,draw=BookTeal,fill=FigureInputFill,line width=1pt},
  gen/.style={draw=FigureMuted,line width=.5pt,-{Stealth[length=1.4mm,width=1mm]}},
  con/.style={draw=FigureMuted,line width=.5pt,dashed,-{Stealth[length=1.4mm,width=1mm]}},
  best/.style={draw=BookTeal,line width=1.2pt,-{Stealth[length=1.8mm,width=1.3mm]}},
  axis label/.style={font=\figurefont\scriptsize,text=FigureMuted}
]
  % 坐标：x=26*t，y=12*g
  \node[cell] (b) at (0,0) {BOS\\$1$};
  \node[cell] (o1) at (26,0) {$\mathtt{o}$\\$0.50$};
  \node[best cell] (a1) at (26,12) {$\mathtt{a}$\\$0.30$};
  \node[cell] (oo2) at (52,0) {$\mathtt{oo}$\\$0.20$};
  \node[cell] (oa2) at (52,12) {$\mathtt{oa}$\\$0.125$};
  \node[best cell] (ab2) at (52,24) {$\mathtt{ab}$\\$0.18$};
  \node[cell] (ooo3) at (78,0) {$\mathtt{ooo}$\\$0.08$};
  \node[cell] (ooa3) at (78,12) {$\mathtt{ooa}$\\$0.05$};
  \node[cell] (oab3) at (78,24) {$\mathtt{oab}$\\$0.075$};
  \node[best cell] (abc3) at (78,36) {$\mathtt{abc}$\\$0.09$};
  \node[cell] (oooo4) at (104,0) {$\mathtt{oooo}$\\$0.032$};
  \node[cell] (oooa4) at (104,12) {$\mathtt{oooa}$\\$0.02$};
  \node[cell] (ooab4) at (104,24) {$\mathtt{ooab}$\\$0.03$};
  \node[cell] (oabc4) at (104,36) {$\mathtt{oabc}$\\$0.0375$};
  \node[best cell] (abcd4) at (104,48) {$\mathtt{abcd}$\\$0.072$};
  \node[done cell] (eos5) at (130,48) {$\mathtt{abcd}\,\mathrm{EOS}$\\$0.0648$};

  % 生成操作（进度不变）
  \foreach \a/\b in {b/o1,o1/oo2,oo2/ooo3,ooo3/oooo4}{\draw[gen] (\a) -- (\b);}
  % 开始或继续约束（进度加一）
  \foreach \a/\b in {o1/oa2,oo2/ooa3,oa2/oab3,ooo3/oooa4,ooa3/ooab4,oab3/oabc4}{\draw[con] (\a) -- (\b);}
  % 交付路径
  \draw[best] (b) -- (a1);
  \draw[best] (a1) -- (ab2);
  \draw[best] (ab2) -- (abc3);
  \draw[best] (abc3) -- (abcd4);
  \draw[best] (abcd4) -- (eos5);

  % 坐标轴说明
  \foreach \t in {0,...,5}{\node[axis label] at ({26*\t},-8) {$t=\t$};}
  \foreach \g in {0,...,4}{\node[axis label,anchor=east] at (-12,{12*\g}) {$g=\g$};}
  \draw[gen] (60,58) -- (70,58) node[right,axis label,text=FigureInk] {生成};
  \draw[con] (84,58) -- (94,58) node[right,axis label,text=FigureInk] {开始或继续约束};
  \draw[best] (122,58) -- (132,58) node[right,axis label,text=FigureInk] {交付路径};
\end{tikzpicture}
